\documentclass[track=junior,level=3,teacher]{kidscode}
%% small helpers for this guide
\newenvironment{guidetable}[1]{\par\noindent\small\renewcommand{\arraystretch}{1.3}%
  \longtable{#1}}{\endlongtable\par}
%% \fitwidth{wide content}: shrink only if it is wider than the line (dyslexia edition)
\newsavebox{\fitbox}
\newcommand{\fitwidth}[1]{\sbox{\fitbox}{#1}\ifdim\wd\fitbox>\linewidth\resizebox{\linewidth}{!}{\usebox{\fitbox}}\else\usebox{\fitbox}\fi}
\begin{document}
\kctitle{Teacher Guide}{Junior Level 3: Bridge to Python. Everything you need before, during and between sessions}
\tableofcontents

%% =====================================================================
\section{Welcome}
Thank you for teaching this course. You do not need to be a programmer or a trained teacher to run
it well. Every session is written out for you: what to print, what to set up, what to say, what to do
when things go wrong, and what to cut if time runs short. If you can read a lesson plan for ten
minutes the day before, you are ready. That is still true in the Python weeks: you will type the
same few short lines the children type, and this guide tells you what every red error message means.

Junior Level 3 is the last junior level, for children aged about 8--9 (US Grade 3) who have finished
Junior Levels 1 and 2 in Scratch. In ten weekly 60-minute sessions they make art with a repeat inside
a repeat, plan programs in parts, collect data and draw bar graphs, play \enquote{teach the robot} to
see how computers learn from examples, and then take their first steps in \textbf{Python}: the
turtle draws the same squares and patterns they know from the Scratch pen, but this time with words.
The level ends with a project, \enquote{My Big Build} (a Scratch game or art app, or a Python
postcard), shown to the family at a \textbf{graduation showcase} in week 10.

\begin{teachernote}[Three things that matter most]
\begin{enumerate}[itemsep=2pt]
  \item \textbf{Bugs are normal, and red text is a clue, not a telling-off.} In the Python weeks the
        computer will complain about a missing letter or bracket many times. Say \enquote{Great, Python
        found a bug for us. Let's read the clue} and mean it. Your calm is the most important thing you teach.
  \item \textbf{Hands off the mouse and the keyboard.} Point, ask and suggest; let the child type,
        even when it is slow. (\enquote{What did you want the turtle to do? What did it do instead?})
  \item \textbf{Change, do not start from nothing.} Every Python program in this level already runs
        and draws something. Children change numbers and color names, or type one short new line.
        That is exactly how real programmers learn a new language.
\end{enumerate}
\end{teachernote}

\begin{teachernote}[The math this level uses]
Everything stays inside Grade 3 math. Numbers are whole numbers up to 1,000 (sizes like 100 or 150
steps). Turns are \textbf{named}, not worked out: a \textbf{quarter turn is 90} and a \textbf{half turn
is 180}. For other shapes children read the number from the turn card (triangle 120, square 90,
hexagon 60). In Scratch, sprites are put in place by dragging them; in Python the turtle walks with
the pen up, using only \texttt{forward}, \texttt{left} and \texttt{right}. There are no x and y
positions and no negative numbers anywhere in this level. The data week uses tally marks and bar
graphs, exactly as in Grade 2--3 math.
\end{teachernote}

\subsection{How the junior track works}
The junior track has three levels of \textbf{10 core weekly sessions}, plus two optional buffer
sessions per level: Junior Level 1 (Little Robots), Junior Level 2 (Game Makers) and Junior Level 3
(Bridge to Python). After that a child moves on to the main track's Level 1. Inside each level the
rhythm is the same: weeks 1--4 new ideas, \textbf{week 5 play day and show-me checklist}, weeks 6--8
new ideas, \textbf{week 9 project studio}, \textbf{week 10 showcase and certificate}. Big ideas come
back again and again (the \enquote{spiral}, see the coverage section): each one is first
\textbf{I}ntroduced, later \textbf{R}evisited and then \textbf{E}xtended.

What makes it junior:
\begin{itemize}[itemsep=2pt]
  \item \textbf{Pages are read aloud by you,} although many Grade 3 children now read parts of them
        themselves. Student pages are picture-first, with about 250 words a week and one short sentence per step.
  \item \textbf{Rose and Pixi videos} (2--3 minutes, one idea each) can open the session, so children
        hear the idea before they need to read anything.
  \item \textbf{Little typing:} in Scratch, short words and nicknames; in Python, numbers, color names
        and at most one short new line per step.
  \item \textbf{No quizzes and no grades.} You watch and check off what you see (show-me checklists).
        Python work is checked by \textbf{matching a picture}: each exercise says what the window should show.
  \item \textbf{Homework is a play-at-home card:} an optional 5--10 minute game with a grown-up.
\end{itemize}

\subsection{How Level 3 bridges Scratch to Python}
Children do not leave Scratch behind: they carry it with them. The bridge is built in three steps.
\begin{enumerate}[itemsep=2pt]
  \item \textbf{Weeks 1--4, still in Scratch.} Nested repeats with the pen and My Blocks
        (the ideas that become \texttt{for} loops and \texttt{def} later), planning in parts,
        collecting data and an unplugged look at how computers learn.
  \item \textbf{Week 5, blocks to words.} A card sort matches each Scratch block with its Python
        twin: \enquote{move 100 steps} and \texttt{t.forward(100)}, \enquote{turn 90 degrees} and
        \texttt{t.left(90)}, \enquote{repeat 4} and \texttt{for i in range(4):}. Children see that
        they already know the ideas; only the spelling is new.
  \item \textbf{Weeks 6--9, Python turtle in Thonny.} Every program already runs. Children press Run,
        compare the window with the picture in the book, change one thing, and press Run again.
        In week 9 each child chooses Scratch or Python for the project.
\end{enumerate}

\begin{center}\small\renewcommand{\arraystretch}{1.35}
\begin{tabularx}{\linewidth}{@{}>{\RaggedRight}p{5.6cm}>{\RaggedRight\arraybackslash}X@{}}
\toprule
\textbf{Scratch block} & \textbf{Python twin}\\
\midrule
when green flag clicked & the green \menu{Run} button in Thonny\\
move 100 steps & \texttt{t.forward(100)}\\
turn left 90 degrees & \texttt{t.left(90)}\\
repeat 4 (with blocks inside) & \texttt{for i in range(4):} with the lines inside pushed in (indented)\\
set pen color & \texttt{t.color("blue")}\\
set pen size to 5 & \texttt{t.pensize(5)}\\
pen up / pen down & \texttt{t.penup()} / \texttt{t.pendown()}\\
a My Block (define \dots) & \texttt{def flower():} (week 8 stretch only)\\
\bottomrule
\end{tabularx}
\end{center}

\begin{guidetable}{@{}>{\bfseries\RaggedRight}p{4.2cm}>{\RaggedRight}p{2.4cm}>{\RaggedRight\arraybackslash}p{\dimexpr\linewidth-6.6cm-4\tabcolsep}@{}}
\toprule
Document & Who & What it is for\\
\midrule\endhead
\bottomrule\endfoot
Teacher Guide (this) & You & Setup, how sessions run, assessment, troubleshooting, pacing.\\
Lesson Plans & You & One 60-minute plan per week (and B1, B2), with words to read aloud.\\
My Coding Book & Children & One short picture chapter per week, read aloud with you.\\
Play-at-Home Cards & Families & One optional 5--10 minute game a week.\\
Answer Key & You only & What each activity should look and sound like; how to use the show-me checklists.\\
Show-Me Checklists & You, per child & Week-1 warm-up, week-5 and week-10 checklists, child self-rating.\\
Project Brief / Project Teacher Guide & Children / you & My Big Build with the plan card and stars / how to run it, the two reference projects, the graduation showcase.\\
Extras & Everyone & Cheat sheet (Scratch and Python side by side), glossary, progress tracker, notes for families, graduation certificate.\\
\end{guidetable}

\textbf{The exercise folders.} Each exercise lives in
\texttt{grade-2-3/\allowbreak{}3-bridge-to-python/\allowbreak{}exercises/\allowbreak{}weekNN/\allowbreak{}ex-name/}.
A \textbf{Scratch} exercise (weeks 1--5) has a \texttt{README.md} (what to build, with the blocks and a
checklist), a \texttt{hints.md} with three hints, and a \texttt{solution.md} for you. A \textbf{Python}
exercise (weeks 6--9) has a \texttt{README.md} with a \enquote{What you should see} section, a
\texttt{hints.md}, the \texttt{starter.py} the child opens and changes, and, for you, a
\texttt{solution.py} and a small test file that the course uses to check the solutions. Children never
run the tests: they compare their window with the README's picture in words. Unplugged games have no
folder: they live in the lesson plan and the workbook. Every PDF also comes in a
\textbf{dyslexia-friendly edition} (clearer font, more spacing, cream paper color).

\subsection{Level 3 goals}
By the end of Junior Level 3 the student can\dots{} (CSTA = CSTA K--12 Computer Science Standards,
revised 2017; band 1B is grades 3--5. AI ideas follow AI4K12's Five Big Ideas in AI. The wording is
abridged.)

\begin{guidetable}{@{}>{\bfseries}p{1.6cm}>{\RaggedRight}p{\dimexpr\linewidth-6.2cm-4\tabcolsep}>{\RaggedRight\arraybackslash}p{4.6cm}@{}}
\toprule
Code & Goal & CSTA / other\\
\midrule\endhead
\bottomrule\endfoot
J3.G1 & use a repeat inside a repeat and My Blocks to make patterns and art & 1B-AP-10, 1B-AP-11\\
J3.G2 & plan a program on paper as parts and steps before building it & 1B-AP-11, 1B-AP-13\\
J3.G3 & collect data, show it in a bar graph and say what it shows & 1B-DA-06, 1B-DA-07\\
J3.G4 & explain that computers learn from examples and can learn unfair mistakes & AI4K12 Big Ideas 3 and 5\\
J3.G5 & read and change short Python turtle programs: numbers, colors and repeat loops & 1B-AP-10, 1B-AP-12\\
J3.G6 & test and debug in blocks and in text, including typos & 1B-AP-15\\
J3.G7 & keep personal information private and give credit for others' work & 1B-NI-05, 1B-AP-14\\
\end{guidetable}

\subsection{The 10-week roadmap}
\begin{guidetable}{@{}>{\bfseries}p{0.9cm}>{\RaggedRight}p{3.4cm}>{\RaggedRight}p{\dimexpr\linewidth-9.9cm-6\tabcolsep}>{\RaggedRight\arraybackslash}p{5.6cm}@{}}
\toprule
Wk & Title & Big idea & Highlights\\
\midrule\endhead
\bottomrule\endfoot
1 & Patterns in Patterns & A repeat inside a repeat makes amazing art. & \enquote{Show me} warm-up (what carried over); clap patterns; The Window and Square Flower with the pen; the one-square flower bug.\\
2 & Plan It on Paper & Big programs are small parts, planned first. & Plan a party in named parts; Party Time built and tested part by part; the silent party bug.\\
3 & Data Detectives & Data is information we collect; a graph tells its story. & Family fruit survey with tally marks and a bar graph; a Scratch survey that counts answers; \enquote{Apple} and \enquote{apple}.\\
4 & How Do Computers Learn? & Computers learn patterns from examples, and copy our mistakes. & \enquote{Teach the robot} with animal cards; a bar graph of its right and wrong guesses; the one-sided robot.\\
5 & Play Day: Blocks to Words & Every Scratch block has a word twin in Python. & Blocks-to-words card sort; bug-hunt relay; my best pattern; \textbf{show-me checklist}.\\
6 & Hello, Python! & Python uses words instead of blocks, and it is picky about spelling. & \enquote{Python is picky}; Thonny and the Run button; Hello, Turtle and Finish the Square; the typo detective.\\
7 & Repeat in Python & \texttt{for} plus an indent is Scratch's repeat block. & \enquote{Indent the line-up}; a square with \texttt{for}; the lonely turn; a Python flower.\\
8 & Colors and My Own Pattern & Change the numbers and colors to make it yours. & Color-by-code cards; color my square; a two-color flower; a misspelled color; my square block with \texttt{def} (stretch).\\
9 & Project Studio & Choose a Scratch game or a Python picture; plan, build, test. & \textbf{Project: M1 plan card, M2 it works}; partner check.\\
10 & Graduation Showcase & Show, explain, celebrate and choose what is next. & Last-minute bug check and typo rescue; \textbf{M3 graduation showcase}; sticker gallery; \textbf{show-me checklist}; certificates.\\
B1 & Robot Olympics & (optional, any time after week 5) & Unplugged sorting races (the line-up race) and maze solving by hand, with a tally scoreboard.\\
B2 & Free Build & (optional, any time after week 5) & Idea cards and a mini plan, then Scratch or Python, the child's choice; swap and test.\\
\end{guidetable}

%% =====================================================================
\section{Before you start: setup checklist}
Allow about an hour, once, a week or more before session 1. You need the internet only to install
Thonny (and Scratch, if it is not on the Mac yet); after that everything works offline. Python is not
needed until week 6, but setting it up now means week 6 starts on time.

\subsection{1. Check Scratch (offline app)}
\begin{checklist}
  \item Scratch 3 should still be on every Mac from Junior Levels 1 and 2. If not: open the
        \textbf{App Store}, search for \textbf{Scratch 3}, check that the maker is the Scratch
        Foundation, and click \menu{Get}. It is free, and children never need a Scratch account.
  \item Open it once and test the Pen: click \menu{Add Extension} (bottom left) and choose \textbf{Pen}.
  \item Keep the \texttt{Scratch Projects} folder in Documents. Saving: \menu{File > Save to your computer}.
        Opening: \menu{File > Load from your computer}.
\end{checklist}

\subsection{2. Install Thonny on every Mac}
Thonny is a free Python editor made for beginners. It includes its own Python, so nothing else is needed.
\begin{checklist}
  \item Go to \texttt{thonny.org} and download the macOS installer (the \texttt{.pkg} file). It is
        free, and no account, email address or sign-in is needed, ever.
  \item Open the \texttt{.pkg} file and click through the installer. Then open Thonny from
        \textbf{Applications}. On the first start, choose your language and the \textbf{Standard}
        initial settings.
  \item \textbf{Make the code bigger:} \menu{View > Increase font size} (or \keys{Cmd}+\keys{+}),
        two or three times, until the Helper can read the code from their chair.
  \item \textbf{Switch to simple mode:} \menu{Tools > Options\dots{} > General}, set \textbf{UI mode}
        to \textbf{simple}, click \menu{OK}, then quit and restart Thonny. Simple mode shows a few big
        buttons (New, Load, Save, Run, Stop and a few more) and hides everything children do not need.
        To change a setting later, click \enquote{Switch to regular mode} at the top right, then come back.
  \item \textbf{Test it offline:} turn Wi-Fi off, click \menu{New} and type these three lines:\par
        \texttt{import turtle}\par\texttt{t = turtle.Turtle()}\par\texttt{t.forward(100)}\par
        Click \menu{Run}, save the file as \texttt{test\_line.py} when asked, and a window opens with a
        line in it. Close that window.
\end{checklist}

\subsection{3. A Coding folder with the student copy}
\begin{checklist}
  \item On every Mac, make a folder called \texttt{Coding} in \textbf{Documents}.
  \item On your own computer, make the \textbf{student copy} of this level: in the course folder run
        \texttt{make student-bundle LEVEL=j3}. It creates
        \texttt{build/\allowbreak student-bundle/\allowbreak kids-coding/}: every exercise folder
        \emph{without} the solutions, plus the project starter \texttt{postcard.py}.
  \item Copy that \texttt{kids-coding} folder (USB stick, AirDrop or a shared drive) into the
        \texttt{Coding} folder on each Mac. Do not give children the whole course folder: it contains
        every answer key, reference solution and teacher guide. Keep that on \textbf{your} Mac.
  \item Check: in Thonny click \menu{Load} and open
        \texttt{Coding/\allowbreak kids-coding/\allowbreak grade-2-3/\allowbreak 3-bridge-to-python/\allowbreak project/\allowbreak starter/\allowbreak postcard.py}.
        Press \menu{Run}: an orange house with a purple roof appears.
\end{checklist}

\subsection{4. A mouse for every Mac}
Plug in a \textbf{mouse} (wired or wireless). Dragging blocks on a trackpad is hard for small hands,
and a mouse makes the Driver and Helper roles clearer: whoever holds the mouse (and, in Python weeks,
sits at the keyboard) is the Driver. For a left-handed child, move the mouse to the left and, if they
like, swap the buttons in \menu{System Settings > Mouse}.

\subsection{5. Print}
\begin{guidetable}{@{}>{\bfseries\RaggedRight}p{3.3cm}>{\RaggedRight\arraybackslash}p{\dimexpr\linewidth-3.3cm-2\tabcolsep}@{}}
\toprule
When & What to print\\
\midrule\endhead
\bottomrule\endfoot
Before week 1 & One per child: My Coding Book (or each chapter the week before), the Cheat Sheet (laminate it if you can: the Scratch-to-Python side is used every Python week), the Progress Tracker, and the Show-Me Checklists (week-1 warm-up page now, the rest later). One per family: Notes for Families. For you: this guide and the Lesson Plans (keep the Answer Key in your folder).\\
Every week & The lesson plan, the workbook chapter, the play-at-home card, and any cards listed under \enquote{Materials} in the plan (clap-pattern cards, survey sheets, animal cards, blocks-to-words cards, spot-the-difference cards, color-by-code cards).\\
Week 5 & The week-5 show-me checklist, one per child.\\
Week 9 & The Project Brief, one per child (the plan card is the page they fill in). The Project Teacher Guide for you.\\
Week 10 & The week-10 show-me checklist, the sticker feedback cards (last page of the Project Teacher Guide) and one graduation certificate per child (write the nickname or first name by hand).\\
\end{guidetable}
Choose the standard or the dyslexia-friendly edition per child. Both have the same pages, so
children can use different editions side by side.

\subsection{6. Room setup}
\begin{itemize}[itemsep=2pt]
  \item \textbf{One Mac per pair}, two chairs side by side, the mouse in the middle. In the Python
        weeks, put the keyboard where the Driver can reach it and the Helper can see the screen.
        Four children need two Macs. A third Mac for your demonstrations is useful but not needed.
  \item \textbf{A timer} (a sand timer or kitchen timer, 6--7 minutes) for Driver and Helper swaps.
  \item \textbf{A clear table for unplugged games}, away from the computers, with space for card sorts
        and a big sheet of paper for the bar graph. Drinks away from the Macs.
  \item \textbf{A floor area} for the clap patterns, \enquote{indent the line-up} (children stand in a
        line, some stepping in) and buffer B1. Masking tape from earlier levels is handy.
  \item Crayons, pencils, sticky notes, scissors, stickers, and squared paper for plans and graphs.
  \item Space to stand up and move for the wiggle break.
\end{itemize}

\subsection{7. Accessibility settings}
\begin{itemize}[itemsep=3pt]
  \item \textbf{Dyslexia-friendly PDFs:} every document has a second edition with the Lexend font,
        extra spacing and a cream background. Offer it to any child; no diagnosis needed.
  \item \textbf{Bigger code in Thonny:} \menu{View > Increase font size} or \keys{Cmd}+\keys{+}.
        Bigger code makes typos much easier to spot.
  \item \textbf{Bigger blocks in Scratch:} the round \enquote{+} zoom button at the bottom right of the code area.
  \item \textbf{A bigger pointer:} \menu{System Settings > Accessibility > Display > Pointer size}.
        Also there: \menu{Color Filters} for children who see colors differently. All course materials
        use color-blind-safe colors and always add a word or an icon to a color. In Python, children
        choose any colors they like; in week 9 they check that everyone can see their picture.
  \item \textbf{Reading and typing support:} read every student page aloud as a habit, and accept
        spoken, pointed or drawn answers everywhere. A child who types slowly can be the Helper who
        spells while you or a partner types, then swap.
\end{itemize}

%% =====================================================================
\section{The 60-minute session}
Every session, including the buffers, has the same five parts. Level 3 sessions are 60 minutes, so
Make it comes in two blocks with the wiggle break in the middle. The activity changes at least every
10--15 minutes.

\begin{timeplan}
  \timeslot{0--10}{Warm-up}{The Rose and Pixi video or an unplugged game (clap patterns, card sorts, acting). Wakes up bodies and brains and plays the new idea before anyone touches a mouse or a key.}
  \timeslot{10--20}{Show me}{The one new idea in kid language, with an analogy, then you build or change the example on the Mac while the children watch, point and predict (\enquote{What will the turtle draw when I press Run?}).}
  \timeslot{20--33}{Make it}{Driver and Helper at the Mac, swapping when the timer runs out (every 6--7 minutes). The first workbook exercise and the bug hunt.}
  \timeslot{33--35}{Wiggle break}{Two minutes away from the screen: stand up, stretch, shake hands, look far away.}
  \timeslot{35--50}{Make it}{The next exercise, then the last few minutes free play: change one number or color and press Run (or click the flag).}
  \timeslot{50--60}{Share}{Each child shows one thing, others say what they like. Exit question, reflection (circle a face), color the progress tracker, hand out the play-at-home card.}
\end{timeplan}

\begin{guidetable}{@{}>{\bfseries}p{2.6cm}>{\RaggedRight\arraybackslash}p{\dimexpr\linewidth-2.6cm-2\tabcolsep}@{}}
\toprule
Part & What it is for\\
\midrule\endhead
\bottomrule\endfoot
Warm-up & Moving and playing first. The unplugged game makes the idea physical (a nested repeat becomes a clap pattern, an indent becomes a step forward), and children who are not yet sure with the keyboard can shine.\\
Show me & One idea only. You model thinking out loud and meeting an error calmly. Children watch with their hands free.\\
Make it & Hands on, in pairs. This is where most learning happens; your job is to ask, not to fix.\\
Wiggle break & Resets attention for the second Make it block and teaches healthy screen habits. Never skip it.\\
Share & Ends on success. Explaining one part of their code out loud is part of the week-10 checklist, and it gives you the best picture of what each child understands.\\
\end{guidetable}

\begin{tip}[Making the 60 minutes work]
\begin{itemize}[itemsep=1pt]
  \item Turn on the Macs and open Scratch or Thonny (with the week's \texttt{starter.py} loaded) \textbf{before} the children arrive.
  \item Protect the Share part. If Make it runs over, cut free play first, never the Share.
  \item Every plan ends with \enquote{If you are short on time}: what to cut without losing the idea.
  \item A tired or wobbly group? Do the wiggle break early. It is never wasted time.
\end{itemize}
\end{tip}

%% =====================================================================
\section{Driver and Helper}
Children work in pairs at one Mac. This doubles the thinking, halves the frustration and builds
the habit of explaining. By Level 3 children know the roles well; the Python weeks give the Helper
a new and very important job.

\begin{pairroles}[The two roles]
\textbf{Driver:} holds the mouse (and, in Python weeks, types), and says out loud what they are doing.\\
\textbf{Helper:} points at the picture or the line in the book and spots bugs. In Scratch weeks the
Helper says which color circle to click next. In \textbf{Python weeks the Helper reads the code
aloud letter by letter}, including the dots and brackets: \enquote{t, dot, f-o-r-w-a-r-d, open bracket,
1-0-0, close bracket}. Python is picky about spelling, so this is real bug-catching work. The Helper
never grabs the mouse or the keyboard: they point and use words.\\
\textbf{Swap} when the timer runs out (every 6--7 minutes), even in the middle of something.
\end{pairroles}

\begin{guidetable}{@{}>{\bfseries}p{2.6cm}>{\RaggedRight\arraybackslash}p{\dimexpr\linewidth-2.6cm-2\tabcolsep}@{}}
\toprule
Children & How to run it\\
\midrule\endhead
\bottomrule\endfoot
1 child & You are the Helper. Resist taking the mouse or the keyboard. Once a session, swap: the child helps while you drive, and you make a small typo for them to catch.\\
2 children & One pair. Seat them so both see the screen equally.\\
3 children & Driver, Helper and \textbf{Checker} (in Scratch, checks each block against the picture; in Python, compares the window with \enquote{What you should see}). Rotate all three jobs on the timer.\\
4 children & Two pairs on two Macs. Visit each other's projects at the end. Change partners every few weeks.\\
Siblings of different ages & They swap on the timer, \textbf{not by age}. The younger child types as long as the older one does. Ask the older child to help with words only: \enquote{Spell it for him, letter by letter, don't type it.}\\
\end{guidetable}

\begin{tip}[Phrases worth teaching]
Helper: \enquote{I'll read it to you, letter by letter}, \enquote{The red text says line 5}, \enquote{Can we press Run to check?}
Driver: \enquote{I'm going to change the 100 to 150}, \enquote{What do you think?}
Both: \enquote{Let's try your idea first, then mine.}
\end{tip}

%% =====================================================================
\section{The threads}
Five kinds of activity run through the level. Each has its own labeled box in the lesson plans.

\begin{unplugged}[No computer]
Clap patterns, planning a party, a family survey, \enquote{teach the robot}, card sorts, spot the
difference and \enquote{indent the line-up}. Unplugged games make the idea physical, give eyes a rest
and let everyone take part. Keep them short (5--10 minutes), give everyone a job, and end with
\enquote{Where will we see this on the computer today?} Buffer B1 is a whole unplugged day.
\end{unplugged}

\begin{bughunt}[Every week]
Children get something that \emph{almost} works: an inner repeat in the wrong place, \enquote{Apple}
and \enquote{apple} counted twice, \texttt{forwad} instead of \texttt{forward}, a line that is not
indented, a color spelled wrong. The steps are on the cheat sheet: \textbf{test} it, \textbf{look}
closely (in Python: read the last red line), \textbf{change one thing}, test again.
Celebrate every catch (\enquote{Bug caught!}).
\end{bughunt}

\begin{freebuild}[Free play]
The last few minutes of Make it are open: change the numbers, the colors, the pen size, the number of
repeats. One rule: change one thing, then press Run (or click the flag). Free play is where children
make a program their own, and it keeps motivation high, so even a pair that is behind gets 3 minutes.
\end{freebuild}

\begin{reflection}[Closing the loop]
Every Share part ends with a short question (\enquote{Which Python word do you remember?}) and a
drawing or a smiley face to circle in the workbook. No writing is needed. Model it yourself:
\enquote{Today I found a bug by reading the red line very slowly.}
\end{reflection}

\begin{safety}[Safety and kindness]
A 2--3 minute talk every week about privacy, data, AI, credit, asking for help and showing work
safely. Make it a conversation, not a lecture. The week-by-week list is in the Digital citizenship
thread section below.
\end{safety}

%% =====================================================================
\section{Assessment plan}
Assessment in the junior track is \textbf{watching, not testing}. There are no quizzes, no
placement tests and no grades. You notice what each child can do while they play, and you use it
to choose what to practice next.

\begin{guidetable}{@{}>{\bfseries\RaggedRight}p{3.6cm}>{\RaggedRight}p{2.4cm}>{\RaggedRight\arraybackslash}p{\dimexpr\linewidth-6.0cm-4\tabcolsep}@{}}
\toprule
What & When & How to use it\\
\midrule\endhead
\bottomrule\endfoot
\enquote{Show me} warm-up & Week 1, 5--10 min & Not graded. Six quick things from Level 2 on the Mac (forever, if touching, a score, a My Block or message, a pen square, a bug fix); two or three per child are enough. Use it to pick pairs and to see which ideas need a reminder in weeks 1--2.\\
Exit question & Every session & In the Share part. One quick question from the lesson plan; listen to the answers.\\
Progress tracker & Every session, 2 min & The child colors the star next to the week's \enquote{I can} statement when they feel ready.\\
Show-me checklist & Week 5 & Six rows. Watch during the play day and check \emph{Not yet}, \emph{Getting there} or \emph{Got it}. Use the results to plan the Python weeks (and B1 if many rows say Not yet).\\
Project stars & Weeks 9--10 & Four rows with 1--3 stars (it works, every must-have, explain one part of the code, planned on the plan card). See the Project Teacher Guide.\\
Show-me checklist & Week 10 & Seven rows, watched during the project and the showcase. Compare with week 5 and keep it for main-track Level 1.\\
Graduation certificate & Week 10 & Every child who finishes the junior track gets one. It celebrates taking part, not a score.\\
\end{guidetable}

\textbf{The three columns} used on every show-me checklist:
\textbf{Not yet} (needs you to show it, or does not try yet),
\textbf{Getting there} (does it with a hint or a reminder),
\textbf{Got it} (does it alone). The Answer Key (\enquote{How to use the show-me checklists})
describes what each looks like for every row and what to do next for \enquote{Not yet}.

\begin{tip}[Valuing the right things]
Praise strategies and effort, not speed: \enquote{You read the red line and found the missing bracket!}
When a child shows you a program, ask \enquote{What does this line do?} before you say
\enquote{Well done}. A child who can explain three lines of Python is further along than one who
copied thirty.
\end{tip}

%% =====================================================================
\section{Differentiation}
Every lesson plan has a support path (needs more time) and an extension path (finished early).
General moves:

\begin{differentiation}
\supportpath Be the Driver yourself while the child points at the next block or reads the next line,
then swap for the last step only. In Python, change only one number per Run, and let the Helper
spell every word. Use the first exercise only (it is always the core one). Repeat the unplugged game.
Zoom the code. Skip the second exercise and come back to it in the next warm-up.
\extensionpath Ask \enquote{Can you change the color? The size? The number of repeats? Add one more
shape?} Let them make a bug for a partner to find. Use the stretch challenge in the plan (in week 8,
\texttt{def}). Let them become the \enquote{bug-hunt helper} who asks questions but never touches
another child's mouse or keyboard.
\end{differentiation}

\begin{guidetable}{@{}>{\bfseries}p{0.9cm}>{\RaggedRight}p{\dimexpr(\linewidth-0.9cm-4\tabcolsep)/2}>{\RaggedRight\arraybackslash}p{\dimexpr(\linewidth-0.9cm-4\tabcolsep)/2}@{}}
\toprule
Wk & \faArrowDown\ Support path & \faArrowUp\ Extension path\\
\midrule\endhead
\bottomrule\endfoot
1 & Start from a square that is already built; add only the outer repeat and its turn. & A pattern with a My Block inside the outer repeat; a hexagon pattern from the turn card.\\
2 & Plan with three named parts only, drawn as pictures. & A plan with a part that is used twice; test each part as soon as it is built.\\
3 & Survey two choices only; build the bar graph with sticky notes. & Ask a second question and compare the two graphs; say what the graph does \emph{not} tell us.\\
4 & Sort with one feature only (legs or no legs). & Find an animal the robot gets wrong and say which example would fix it.\\
5 & Play day with help cards; show-me rows watched one at a time; sort five block-and-word pairs only. & Be the bug-hunt helper in the relay; write the Python twin for a new block.\\
6 & You type, the child spells; change one number per Run. & Draw a letter L, then a whole square without a loop (the long way, to feel why loops help).\\
7 & The loop is already there; the child changes the number in \texttt{range} and the side length. & A triangle and a hexagon from the turn card; a loop inside a loop (light).\\
8 & Change the color names only, from the color-by-code cards. & \texttt{def} for a repeated shape, used two or three times.\\
9 & Core path: start from a template (a Scratch game set up by you, or \texttt{postcard.py}); one must-have at a time. & Stretch path: a \texttt{def} for a repeated shape (or a My Block in Scratch).\\
10 & Practice the show-and-tell with you first; point and say one sentence. & Tell the story of a bug they fixed; help another child set up.\\
\end{guidetable}

%% =====================================================================
\section{Digital citizenship thread}
Every week has a short safety, kindness or wellbeing moment (a \enquote{Digital citizenship \&
safety} box in the lesson plan). Keep it to 2--3 minutes and make it a conversation.

\begin{guidetable}{@{}>{\bfseries}p{0.9cm}>{\RaggedRight}p{4.4cm}>{\RaggedRight\arraybackslash}p{\dimexpr\linewidth-5.3cm-4\tabcolsep}@{}}
\toprule
Wk & Theme & The key message\\
\midrule\endhead
\bottomrule\endfoot
1 & My own pattern, not a copy & Getting ideas from others is fine; make it your own and say where the idea came from.\\
2 & Teamwork & Everyone gets a job, and every job matters.\\
3 & Only ask what you need & A survey asks only what it needs (favorite fruit), never personal information such as a full name, address, school or birthday.\\
4 & AI can be wrong, and it is not a friend & A computer that learned from examples can make mistakes, even unfair ones. A chatbot is not a friend: never tell it your name, address or school.\\
5 & Proud of my work & What am I proud of so far? Saying it out loud is a skill too.\\
6 & Saving files with clear names & Name files so you can find them (\texttt{tiger\_square.py}), with a nickname, never a full name.\\
7 & Asking for help well & Say what you expected, what happened instead, and what you already tried.\\
8 & Credit for remixes & Whose idea did I remix? Say so: \enquote{I started from Zoe's flower and changed the colors.}\\
9 & Can everyone see it? & Big shapes, thick lines, and colors plus shapes or words, so everyone can enjoy it.\\
10 & Next steps safely & If a website or app asks for your name, age, school or a photo, stop and ask a grown-up.\\
\end{guidetable}

%% =====================================================================
\section{Troubleshooting}
\begin{guidetable}{@{}>{\RaggedRight}p{5.4cm}>{\RaggedRight\arraybackslash}p{\dimexpr\linewidth-5.4cm-2\tabcolsep}@{}}
\toprule
\textbf{Problem} & \textbf{What to do}\\
\midrule\endhead
\bottomrule\endfoot
\multicolumn{2}{@{}l}{\textbf{\color{kcaccent}Scratch (weeks 1--5 and the Scratch project)}}\\ \midrule
A block will not snap on. & Drag it by its left end and move it slowly under the last block. Wait for the gray shadow, then let go. A hat block (round top) only goes at the top.\\
Nothing happens when the green flag is clicked. & The script needs \enquote{when green flag clicked} on top, and every block must be snapped on. Check the script is under the right sprite.\\
The pattern is drawn in one place only, or once. & The inner repeat is in the wrong place, or the turn between the shapes is outside the outer repeat (the week-1 bug).\\
The pen blocks are missing. & Click \menu{Add Extension} (bottom left), then \textbf{Pen}. Each new project needs it again.\\
The cat is lost off the stage. & A small piece always stays at the edge: find it and drag it back to the middle.\\
A variable does not go back to 0. & Add \enquote{set score to 0} under \enquote{when green flag clicked}.\\
\midrule
\multicolumn{2}{@{}l}{\textbf{\color{kcaccent}Thonny and Python (weeks 6--9 and the Python project)}}\\ \midrule
Nothing happens when we press Run. & Check the green \menu{Run} button was pressed (or \keys{F5}). A new file must be saved before it runs: when the Save window appears, save it in the \texttt{Coding} folder with a clear name. If Run is gray, a program is still running: press the red \menu{Stop}, then Run again.\\
Red text appears in the Shell at the bottom. & That is Python telling us about a bug. Read the \textbf{last} red line aloud together: it names the problem, and the line above it says which line number. Look at that line \emph{and the line above it}. In regular mode, \menu{View > Assistant} explains errors in friendlier words.\\
The red text says \enquote{Did you mean: 'forward'?} & A spelling mistake: \texttt{t.forwad} instead of \texttt{t.forward}. Python guesses the word you meant. Fix the spelling and Run again.\\
\texttt{NameError: name 'T' is not defined} & Capitals matter: \texttt{t} and \texttt{T} are different names. Python words are lower case (\texttt{t.forward}, \texttt{t.left}).\\
\texttt{SyntaxError: '(' was never closed} & A bracket is missing. Count them with the Helper: every \texttt{(} needs a \texttt{)}, and every \texttt{"} needs a partner.\\
\texttt{IndentationError: expected an indented block} & The line under \texttt{for i in range(4):} is not pushed in. Click at the start of that line and press \keys{Tab} (Thonny adds 4 spaces).\\
\texttt{IndentationError: unexpected indent} (or \enquote{unindent does not match}) & A line is pushed in when it should not be, or by a different number of spaces. Lines inside the loop need exactly 4 spaces; other lines none. Select the line and press \keys{Shift}+\keys{Tab} to move it back.\\
The square is a straight line, or only one side is drawn. & \texttt{t.left(90)} is not pushed in, so it is outside the loop (the week-7 bug). Push it in to line up with \texttt{t.forward}.\\
\texttt{TurtleGraphicsError: bad color string: bleu} & A color name is spelled wrong (the week-8 bug). Use the color names on the cheat sheet: orange, blue, purple, green, gold, black.\\
\texttt{AttributeError:\allowbreak{} module 'turtle' has no attribute 'Turtle'} & The file was saved as \texttt{turtle.py}. Rename it (for example \texttt{my\_square.py}) and delete any \texttt{turtle.py} in that folder.\\
We cannot see the turtle window. & It is hiding behind Thonny. Press \keys{Cmd}+\keys{Tab}, or click the Python rocket icon in the Dock.\\
Pressing Run again does nothing new, or two windows appear. & Close the turtle window first (the red dot at its top left), then press Run. If it will not close, press \menu{Stop} in Thonny.\\
The window opens and closes at once. & The last line \texttt{turtle.done()} is missing. Add it back at the very end, not pushed in.\\
The menus are gone. & That is simple mode. Click \enquote{Switch to regular mode} at the top right of Thonny when you (not the child) need a menu, and switch back afterwards.\\
\midrule
\multicolumn{2}{@{}l}{\textbf{\color{kcaccent}Saving and opening}}\\ \midrule
Where do I save? & Scratch: \menu{File > Save to your computer}, in \texttt{Scratch Projects}. Python: the \texttt{starter.py} in the exercise folder is saved where it is (\menu{Save}). New files go in \texttt{Coding}. Use a nickname and the project in the file name (\texttt{tiger\_postcard.py}), never a real full name.\\
I cannot find a program. & Scratch: \menu{File > Load from your computer}. Thonny: \menu{Load}, then look in \texttt{Coding}. Also check Downloads. Save at the end of every session, before Share.\\
\midrule
\multicolumn{2}{@{}l}{\textbf{\color{kcaccent}The Mac}}\\ \midrule
An app asks to use the microphone or camera. & Click \menu{Don't Allow}; we do not need them.\\
Screen Time blocks Scratch or Thonny. & Add them to \enquote{Always Allowed} in \menu{System Settings > Screen Time}, or schedule sessions outside downtime.\\
Typing is very slow. & Normal at this age. The Helper spells, the Driver types, and nobody hurries. Starters keep typing short on purpose.\\
\end{guidetable}

%% =====================================================================
\section{Pacing, buffers and holidays}
Junior Level 3 has \textbf{10 core sessions}, one a week. Weeks with no session (holidays, illness)
do not count. The buffers \textbf{B1} (Robot Olympics, unplugged sorting races and maze solving) and
\textbf{B2} (Free Build, Scratch or Python) are \textbf{optional} and can go anywhere after week 5: use
B1 when a group needs more unplugged practice or after a long break, and B2 when children want more
free building before the project.

\begin{tip}[Rules of thumb]
\begin{itemize}[itemsep=2pt]
  \item \textbf{Never skip weeks 5 and 10.} They hold the show-me checklists, and week 10 is the
        graduation showcase every child has been working towards.
  \item \textbf{Do not split week 6 and week 7 with a long break} if you can help it: the first
        Python steps stick best when they come close together.
  \item \textbf{Lose a session?} Drop a buffer if you planned one. Otherwise do the warm-up of the
        missed week as the next warm-up and move on: the ideas come back later.
  \item \textbf{One child missed a session?} Read the workbook chapter together for 5 minutes and pair
        them as Helper first next time.
  \item \textbf{After a long break} start with 5 minutes of changing numbers in a program that already
        runs, before anything new.
\end{itemize}
\end{tip}

\subsection{One school year or two}
\begin{itemize}[itemsep=2pt]
  \item \textbf{One school year:} all three junior levels (30--36 sessions at one a week), using the
        buffers for holidays and \enquote{we need another week} moments.
  \item \textbf{Two school years} (recommended when a child starts in early Grade 2):
        Grade 2 = Junior Levels 1 and 2 with both buffers; Grade 3 = Junior Level 3, then main-track
        Level 1. The calendar below fits this plan: Level 3 in the fall of Grade 3, main-track Level 1
        after the winter break.
\end{itemize}

\subsection{Sample calendar for fall 2027}
The example uses \textbf{Saturday mornings}, starting on 11 September 2027, the Saturday after
Labor Day weekend (Labor Day is Monday 6 September, so its Saturday, 4 September, comes before the
start). The 10-week plan fits before Thanksgiving with no skips. With both buffers it runs into
Thanksgiving weekend, so it skips the Saturday after Thanksgiving (27 November). If you meet on a
weekday, shift each date to your day and skip the holidays that fall on it.

\begin{center}\small\renewcommand{\arraystretch}{1.2}
\fitwidth{\begin{tabular}{@{}>{\bfseries}l l l@{}}
\toprule
Saturday & 10-week plan & With B1 and B2 (12 sessions)\\
\midrule
Sep 11 & W1 Patterns in Patterns & W1\\
Sep 18 & W2 Plan It on Paper & W2\\
Sep 25 & W3 Data Detectives & W3\\
Oct 2 & W4 How Do Computers Learn? & W4\\
Oct 9 & W5 Play Day: Blocks to Words & W5\\
Oct 16 & W6 Hello, Python! & B1 Robot Olympics\\
Oct 23 & W7 Repeat in Python & W6\\
Oct 30 & W8 Colors and My Own Pattern & W7\\
Nov 6 & W9 Project Studio & W8\\
Nov 13 & W10 Graduation Showcase & B2 Free Build\\
Nov 20 & (next: main-track Level 1) & W9 Project Studio\\
Nov 27 & \emph{skip: Thanksgiving weekend} & \emph{skip: Thanksgiving weekend}\\
Dec 4 & & W10 Graduation Showcase\\
\bottomrule
\end{tabular}}
\end{center}
\begin{itemize}[itemsep=2pt]
  \item \textbf{10-week plan:} 10 sessions from Sep 11 to Nov 13. Indigenous Peoples' Day / Columbus
        Day (Mon Oct 11) and Veterans Day (Thu Nov 11) do not fall on a Saturday.
  \item \textbf{With both buffers:} 12 sessions from Sep 11 to Dec 4. B1 right after the week-5
        checklist lets you practice the rows marked \enquote{Not yet}; B2 right before the project gives
        children a free build to warm up their ideas. Because Thanksgiving falls between weeks 9 and 10,
        start week 10 with 5 minutes of reopening and testing the project.
\end{itemize}

%% =====================================================================
\section{Coverage: what is taught when}

\subsection{(a) Goals across the weeks}
\cvI\ = introduced,\quad \cvR\ = revisited,\quad \cvE\ = extended. Every goal is introduced and
practiced before the week-10 showcase; the project and the showcase bring them together.

\begin{center}\small\renewcommand{\arraystretch}{1.35}
\fitwidth{\begin{tabular}{@{}>{\bfseries}l*{10}{c}@{}}
\toprule
Goal & W1 & W2 & W3 & W4 & W5 & W6 & W7 & W8 & W9 & W10\\
\midrule
G1 Nested repeat, My Blocks & \cvI &  &  &  & \cvR &  & \cvE & \cvE & \cvR & \\
G2 Plan in parts &  & \cvI &  &  & \cvR &  &  &  & \cvE & \cvR\\
G3 Data and bar graphs &  &  & \cvI & \cvR & \cvR &  &  &  &  & \\
G4 How computers learn &  &  & \cvI & \cvE & \cvR &  &  &  &  & \\
G5 Read and change Python &  &  &  &  & \cvI & \cvR & \cvE & \cvE & \cvR & \cvR\\
G6 Test and debug & \cvI & \cvR & \cvR & \cvR & \cvE & \cvE & \cvR & \cvR & \cvR & \cvR\\
G7 Privacy and credit & \cvI &  & \cvE & \cvE & \cvR & \cvR & \cvR & \cvR & \cvR & \cvR\\
\bottomrule
\end{tabular}}
\end{center}
{\small Notes: G1 starts in Scratch in week 1 and is extended in Python: a nested loop in the
Python flower (week 7) and \texttt{def} as a My Block (week 8, stretch). G2 is named in week 2,
practiced in the week-5 pattern plan and extended with the plan card in week 9. G3 and G4 are compact
strands: G3 is the week-3 survey and the week-4 bar graph of the robot's guesses; G4 starts with an
unfair, one-sided survey in week 3 and becomes \enquote{teach the robot} in week 4. Both are on the
week-5 checklist, and buffer B1 adds a tally scoreboard. G5 starts with the blocks-to-words card sort
in week 5. G6 is extended twice: the bug-hunt relay in week 5 and typos in text in week 6. G7 is
extended in week 3 (only ask what you need) and week 4 (AI can be wrong; never tell a chatbot personal
information). The generated table in (b) shows exactly which lesson objectives serve each goal.\par}

\subsection{(b) Objectives and where they are taught}
The table below is generated automatically from the lesson plans (every \enquote{The student
can\dots} objective with its goal code and week).

\kcinput{generated/coverage}

\subsection{(c) The spiral map: how big ideas come back}
Junior Level 3 extends what Junior Levels 1 and 2 began and hands each idea on to the main track.

\begin{guidetable}{@{}>{\bfseries\RaggedRight}p{2.8cm}>{\RaggedRight}p{\dimexpr(\linewidth-2.8cm-8\tabcolsep)/4}>{\RaggedRight}p{\dimexpr(\linewidth-2.8cm-8\tabcolsep)/4}>{\RaggedRight}p{\dimexpr(\linewidth-2.8cm-8\tabcolsep)/4}>{\RaggedRight\arraybackslash}p{\dimexpr(\linewidth-2.8cm-8\tabcolsep)/4}@{}}
\toprule
Big idea & Junior L1 & Junior L2 & Junior L3 & Main track\\
\midrule\endhead
\bottomrule\endfoot
Sequence and algorithms & W1 \cvI\ (unplugged), W2 \cvR & W8 \cvR\ (game states) & W2 \cvE\ (planning in parts) & L1 W1, L2 W10\\
Loops & W4 \cvI\ (repeat), W8 \cvR\ (pen) & W1 \cvE\ (forever) & W1 \cvE\ (nested), W7 \cvR\ (for) & L1 W2, L1 W8\\
Events & W2 \cvI\ (green flag), W6 \cvE\ (click, key) & W6 \cvE\ (messages) & W6--8 \cvR\ (Run) & L1 W3\\
Conditionals & --- & W2 \cvI\ (if touching), W8 \cvR & W4 \cvR\ (sorting rules) & L1 W3, L1 W9\\
Variables and data & B1 (numbers stand for colors) & W3 \cvI\ (score), W4 \cvI\ (random) & W3 \cvE\ (survey data) & L1 W4, L2 W3\\
Abstraction & W9 storyboard & W7 \cvI\ (My Blocks) & W1 \cvR, W8 def (stretch) & L1 W5, L1 W10\\
Debugging & every week, W3 \cvI & W5 relay, playtesting & W6 \cvE\ (typos in text) & every week\\
AI & --- & --- & W4 \cvI\ (unplugged) & L3 W9\\
Citizenship & nicknames, breaks, credit, passwords & fairness, kindness, credit, accessibility & privacy, data, AI safety, credit & every week\\
\end{guidetable}

%% =====================================================================
\section{Next steps: into main-track Level 1}
A child who finishes Junior Level 3 is ready for the main track's \textbf{Level 1}. Scratch, events,
loops, variables and a first look at Python are already familiar, so the start will feel like
meeting old friends.

\begin{itemize}[itemsep=2pt]
  \item \textbf{Review weeks 1--5 quickly.} Use main-track Level 1 weeks 1--5 as a review: two weeks per
        session, or as homework. Then teach weeks 6--12 in full.
  \item \textbf{Expect more reading and math.} The main track assumes Grade 4--5 reading and math. A
        child still in Grade 3 may need you to read aloud and to help with angles (the \enquote{Math
        helper} boxes) when they get there. Main-track pages also use x and y positions and negative
        numbers, which are new.
  \item \textbf{Keep the week-10 show-me checklist.} Rows marked \enquote{Not yet} tell you what to
        practice first, and the Python rows show how much of Thonny is already familiar.
  \item \textbf{Keep the projects and Thonny.} The \texttt{Coding} and \texttt{Scratch Projects} folders
        come along, so children can open and remix their Big Build in free play. In the main track,
        children switch Thonny back to regular mode when the course says so.
  \item \textbf{A long break first?} Start with buffer B2 of this level (Free Build) before the first
        new idea.
\end{itemize}

\end{document}
